\introduction

\subsection*{Актуальность темы}

Уведомления в прикладных системах используются для подтверждения операций, передачи изменений статусов, напоминаний и транзакционных сообщений. Для таких сценариев важен не только вызов внешнего канала связи, но и сохранность принятого события, повторяемость обработки и восстановление итогового статуса. При наличии нескольких интеграций и каналов уведомление становится распределённым процессом: команда принимается, событие фиксируется, доставка запускается, канальный сервис выполняет отправку, а результат сохраняется в истории [1; 2].

Практическая сложность такого процесса связана с тем, что отказ может возникнуть на разных участках одного маршрута. Начальный контур может сохранить событие, но завершиться до публикации сообщения в Kafka. Брокер может повторно доставить одно и то же сообщение сервису канальной отправки. Внешний канальный провайдер может вернуть ошибку или тайм-аут уже после фактической отправки сообщения. Поэтому платформа должна не только передавать сообщение в канал, но и фиксировать принятые события, подавлять повторную обработку там, где это возможно, выполнять контролируемые повторные попытки, поддерживать fallback-маршрутизацию и сохранять итоговый статус.

При прямой цепочке синхронных вызовов сбой канального сервиса или временная недоступность зависимости напрямую отражаются на клиентском запросе. Если каждая прикладная система самостоятельно интегрируется с внешними канальными провайдерами, то правила доставки, проверки профиля получателя, отмены и повторной обработки начинают дублироваться в разных частях корпоративного контура. Это приводит к расхождению поведения между сервисами и каналами связи, усложняется сопровождение и затрудняетс расследование инцидентов.

Для платформы уведомлений недостаточно применять механизмы надёжности отдельно друг от друга. Их необходимо связать с предметной моделью уведомительного процесса, канальной логикой, проверкой профиля получателя и историей статусов. Именно эта связка определяет актуальность разработки распределённой платформы, в которой приём команды, запуск доставки, канальная обработка и фиксация результата разделены между сервисными контурами, но остаются частью единого управляемого процесса.

\subsection*{Актуальность темы}

Уведомления в прикладных системах используются для подтверждения операций, передачи изменений статусов, напоминаний и транзакционных сообщений. Для таких сценариев важен не только вызов внешнего канала, но и сохранность принятого события, повторяемость обработки и восстановление итогового статуса. При наличии нескольких интеграций и каналов уведомление становится распределённым процессом: команда принимается, событие фиксируется, доставка запускается, канальный сервис выполняет отправку, а результат сохраняется в истории \cite{kleppmann_ddia,newman_microservices}.

Практическая сложность связана с разными типами сбоев на одном маршруте. Начальный контур может сохранить событие, но завершиться до публикации сообщения в Kafka. Брокер может повторно доставить одно и то же сообщение сервису канальной отправки. Внешний провайдер может вернуть timeout уже после фактической отправки. Поэтому платформа должна фиксировать принятые события, подавлять повторную обработку там, где это возможно, выполнять контролируемые повторы, поддерживать fallback-маршрутизацию и сохранять итоговый статус. При прямой цепочке синхронных вызовов сбой канального сервиса или временная недоступность зависимости напрямую отражаются на клиентском запросе.

Для платформы уведомлений недостаточно применять эти механизмы отдельно. Их нужно связать с предметной моделью уведомительного процесса, канальной логикой, проверкой профиля получателя и историей статусов. Эта связка определяет архитектуру разрабатываемой системы.

\subsection*{Объект, предмет, цель и задачи исследования}

\textbf{Объект исследования} -- процессы формирования, маршрутизации, доставки и учёта пользовательских уведомлений в распределённой информационной системе.

\textbf{Предмет исследования} -- архитектурные и программные механизмы, обеспечивающие сохранность уведомительных событий, управляемую многоканальную обработку и фиксацию результатов доставки.

\textbf{Цель работы} -- разработать и реализовать архитектуру распределённой платформы уведомлений, обеспечивающую устойчивую обработку уведомительных событий, контролируемую канальную доставку и прослеживаемость результатов.

Для достижения поставленной цели определены следующие задачи:
\begin{enumerate}
    \item проанализировать уведомление как распределённый процесс и сформулировать границы ответственности платформы;
    \item определить требования к сохранности принятого события, повторной обработке, мультиканальности и истории статусов;
    \item сравнить архитектурные варианты и выбрать схему межсервисного взаимодействия для командного и фонового контуров;
    \item спроектировать доменную модель уведомительного процесса и разделение хранилищ по сервисным контурам;
    \item реализовать сервисы в рамках выбранного архитектурного подхода;
    \item реализовать механизмы сохранности и надёжной обработки;
    \item реализовать два независимых канальных контура: email и SMS;
    \item реализовать запись и чтение истории доставки;
    \item провести стендовую и нагрузочную проверку полного маршрута обработки.
\end{enumerate}

\subsection*{Практическая значимость}

Практическая значимость работы состоит в реализации единого контура обработки уведомительных событий для нескольких каналов связи. Платформа централизует приём команд, запуск доставки, проверку профиля получателя, канальную обработку и сохранение истории статусов. За счёт этого прикладным системам не требуется самостоятельно реализовывать интеграции с каждым канальным провайдером и повторять логику обработки ошибок, повторных попыток и дедупликации.

Мультиканальность решения выражается в том, что разные каналы доставки подключаются к общей модели уведомительного события, запуска доставки и истории, но обрабатываются независимыми канальными сервисами. В текущей реализации поддержаны email и SMS провайдеры, однако архитектура системы не привязывает жизненный цикл уведомления к конкретному способу доставки. Это позволяет добавлять новые каналы связи через расширение канального контура и шаблонного представления без переработки основного процесса обработки события что делает систему более гибкой и масштабируемой.

Разработанное решение снижает риск потери принятого события внутри платформы, поскольку переход от сохранения события к асинхронной обработке выполняется через механизм надёжной публикации. Повторная обработка сообщений допускается архитектурой, но ограничивается ключами идемпотентности, дедупликацией через Redis и уникальными ограничениями в хранилищах. Такой подход позволяет использовать модель обработки не менее одного раза без неконтролируемого роста дублей.

Отдельную практическую ценность имеет единая история доставки. Запись статусов позволяет восстановить маршрут обработки от входной команды до итогового результата, определить стадию сбоя и отделить временные ошибки от окончательных. Это уменьшает трудоёмкость эксплуатационной диагностики и делает поведение платформы проверяемым не только по логам отдельных сервисов, но и по накопленной истории событий.

\subsection*{Границы работы}

В качестве используемых каналов связи выбраны email и SMS. Гарантия доставки ограничена внутренним контуром платформы и относится к сохранности принятого события, безопасной повторной обработке и фиксации статуса. Физическая доставка конечному пользователю зависит от внешнего провайдера и подтверждается только канальным статусом. Вопросы production-развёртывания, полного промышленного контура безопасности и подключения дополнительных каналов связи относятся к дальнейшему развитию.
